\section{Conformal Proofs}\label{sec:CP}

\subsection{Proof of \texorpdfstring{\theoremref{theorem:conformal_method}}{Conformal Probabilistic Safety Verification Theorem}}\label{apd:conformal_theorem_proof}

\begin{proof}

\theoremref{theorem:conformal_method} is a straightforward application of the split conformal prediction method detailed in \citet{MAL-101}, where we set the conformal ``input'' $x=x$, ``output'' $y=-\costFunction(x,0)$, ``score function'' $s(x,y)=y=-\costFunction(x,0)$, ``size of the calibration set'' $n=N$, and ``user-chosen error rate'' $\alpha=\frac{k+1}{N+1}$.
The conformal $\hat{q}$ is then computed as the $\frac{\lceil(N+1)(1-\alpha)\rceil}{N}$ quantile of the calibration scores $-\costFunction(x_{1:N},0)$.
The quantile is $\frac{\lceil(N+1)(1-\alpha)\rceil}{N} = \frac{\lceil(N+1)(1-\frac{k+1}{N+1})\rceil}{N} = \frac{\lceil(N+1)(\frac{N-k}{N+1})\rceil}{N} = \frac{N-k}{N}$, where we have defined $k$ in the procedures in \sectionref{sec:scenario-based_method} as the number of scores $-\costFunction(x_i,0) \ge 0$.
Thus, this quantile corresponds precisely to the largest \textit{negative} score, so we know that $\hat{q} < 0$.
Theorem 1 in \citet{MAL-101} then yields:
\begin{align}
    \underset{ \left( x_{1:N},x \right) \in \safeset }{\mathbb{P}} \left( -\costFunction(x,0) \le \hat{q}  \right) &\ge 1-\alpha \nonumber \\
    \underset{ \left( x_{1:N},x \right) \in \safeset }{\mathbb{P}} \left( \costFunction(x,0) \ge -\hat{q}  \right) &\ge 1-\frac{k+1}{N+1} \nonumber \\
    \underset{ \left( x_{1:N},x \right) \in \safeset }{\mathbb{P}} \left( \costFunction(x,0) > 0  \right) &\ge \frac{N-k}{N+1} \label{eq:cp_proof_coverage}
\end{align}
where Equation \eqref{eq:cp_proof_coverage} follows from the line preceding it because if $\costFunction(x,0) \ge -\hat{q}$, then certainly $\costFunction(x,0) > 0$.
This coverage property result is precisely the same as described in \remarkref{remark:CP_coverage}.
Furthermore, Section 3.2 in \citet{MAL-101} yields:
\begin{align}
    \underset{ x \in \safeset }{\mathbb{P}} \left( \costFunction(x,0) > 0  \right) &\sim \text{Beta}(N+1-l,l),\quad l= \lfloor (N+1)\alpha\rfloor \nonumber \\
    \underset{ x \in \safeset }{\mathbb{P}} \left( \costFunction(x,0) > 0  \right) &\sim \text{Beta}(N-k,k+1)
\end{align}
which is precisely the result of \theoremref{theorem:conformal_method}.
\end{proof}

\subsection{Proof that Split Conformal Prediction Reduces to Robust Scenario Optimization}\label{apd:reduction_proof}
Here, we show that split conformal prediction, in full generality, reduces to a robust scenario optimization problem.
We hope that this insight will encourage future research on the close relationship between the highly related but disparate fields of conformal prediction and scenario optimization.
In split conformal prediction, we first define a score function $s(x,y) \in \mathbb{R}$ which is meant to reflect the uncertainty for a model input $x$ and corresponding model output $y$.
Then, we sample an i.i.d. calibration set $(X_1,Y_1),...,(X_n,Y_n)$ and compute $\hat{q}$ as the $\frac{\lceil(n+1)(1-\alpha)\rceil}{n}$ quantile of the calibration scores $s(X_1,Y_1),...,s(X_n,Y_n)$, where $\alpha \in [0, 1]$ is a user-chosen error rate.
For a new i.i.d. sample $X_\text{test}$, we construct a prediction set $C(X_\text{test})=\{y: s(X_\text{test},y) \le \hat{q}\}$.
Theorem 1 in \citet{MAL-101} provides the following coverage property: $\mathbb{P} \left( Y_{\text{test}} \in C(X_\text{test}) \right) \ge 1-\alpha$.
This follows from the more powerful property, first introduced in \citet{pmlr-v25-vovk12}, which we prove reduces to a robust scenario-based result after:
\begin{equation}\label{eq:conditional_coverage}
    \mathbb{P} \left( Y_\text{test} \in C(X_\text{test}) | \{ (X_i, Y_i) \}^n_{i=1} \right) \sim \text{Beta}(n+1-l, l), \quad l=\lfloor (n+1)\alpha \rfloor
\end{equation}

\begin{proof}
    
To show that split conformal prediction reduces to a scenario-based approach, consider the CCP \eqref{eq:CCP} in \lemmaref{lemma:CCP} in \appendixref{apd:SO}, where $h=(x,y)$ and $f(h)=f\left((x,y)\right)=s(x,y)$.
That is, we want to find a probabilistic upper-bound on samples of the score function.
Then, consider the corresponding SP \eqref{eq:SP} in \lemmaref{lemma:CCP} where the $n$ sampled calibration scores $s(X_1, Y_1), ..., s(X_n, Y_n)$ forms our set of constraints.
Remove the $k=\lfloor (n+1)\alpha-1 \rfloor$ largest scores, where $\alpha$ is the user-chosen error rate in split conformal prediction.
The largest remaining score will be the $\frac{n-k}{n}$ quantile.
$\frac{n-k}{n}=\frac{n-\lfloor (n+1)\alpha-1 \rfloor}{n}=\frac{n+\lceil -((n+1)\alpha-1) \rceil}{n}=\frac{\lceil n -(n+1)\alpha + 1 \rceil}{n}=\frac{\lceil (n+1)(1-\alpha) \rceil}{n}$, which is precisely the same quantile as $\hat{q}$ in split conformal prediction. Thus, the solution to the SP is $g^*_{N,k}=\hat{q}$. \lemmaref{lemma:CCP} tells us that for a violation parameter $\epsilon \in (0, 1)$ and a confidence parameter $\beta \in (0, 1)$ that satisfies the relationship in Equation \eqref{eq:parameters_relationship}, with probability at least $1-\beta$, the following holds:
\begin{equation}\label{eq:scenario_based_conformal_prediction}
    \mathbb{P}\left( Y_\text{test} \in C(X_\text{test}) | \{ (X_i, Y_i) \}^n_{i=1} \right) \ge 1-\epsilon
\end{equation}

This is equivalent to Equation \eqref{eq:conditional_coverage}. To see why, note that the cumulative distribution function of the Beta distribution in Equation \eqref{eq:conditional_coverage} is given in terms of $k$ by the incomplete beta function ratio $I_x(n-k, k+1)=\sum^n_{j=n-k}\binom{n}{j}x^j(1-x)^{n-j}$ from \citet[\href{https://dlmf.nist.gov/8.17.E5}{(8.17.5)}]{NIST:DLMF}.
Changing the index $i=n-j$ yields $I_x(n-k, k+1)=\sum^k_{i=0}\binom{n}{n-i}x^{n-i}(1-x)^{n-(n-i)}=\sum^k_{i=0}\binom{n}{i}x^{n-i}(1-x)^{i}$.
Thus, Equation \eqref{eq:conditional_coverage} is equivalent to the claim that for any violation parameter $\epsilon \in (0, 1)$ and confidence parameter $\beta \in (0, 1)$, $\mathbb{P}\left( Y_\text{test} \in C(X_\text{test}) | \{ (X_i, Y_i) \}^n_{i=1} \right) \ge 1-\epsilon$ (Equation \eqref{eq:scenario_based_conformal_prediction}) holds with probability at least $1-\beta$ as long as $\beta \ge I_{1-\epsilon}(n-k, k+1)=\sum^k_{i=0}\binom{n}{i}\epsilon^{i}(1-\epsilon)^{n-i}$ (Equation \eqref{eq:parameters_relationship}).

\end{proof}

\subsection{Proof of \texorpdfstring{\lemmaref{lemma:conformal_method}}{Conformal Probabilistic Safety Verification Lemma}}\label{apd:conformal_lemma_proof}

\begin{proof}
    The conformal probabilistic safety verification method in \sectionref{sec:conformal_method} is nothing more than a specific formulation of the general split conformal prediction method in \appendixref{apd:reduction_proof}, which we have proven provides a result that is equivalent to the robust scenario optimization result in \lemmaref{lemma:CCP}.
    The result in \lemmaref{lemma:CCP}, when formulated in the context of the conformal probabilistic safety verification method in \sectionref{sec:conformal_method} and noting that $\forall (\state,\tvar), \costFunction(\state,\tvar) \le V(\state, \tvar)$ from Equation \eqref{eq:valuefunc}, is precisely \lemmaref{lemma:conformal_method}.
\end{proof}